\documentclass[UTF8,a4paper,12pt]{ctexart}
\usepackage{geometry,amsmath,booktabs,hyperref}
\geometry{left=2.5cm,right=2.5cm,top=2.5cm,bottom=2.5cm}
\title{高性能芯片热管理系统问题四：权重敏感性与鲁棒设计}
\author{数学建模项目组}\date{\today}
\begin{document}\maketitle

\section{问题描述}
问题二得到三个代理模型 $\hat R(x)$、$\hat P(x)$ 和 $\hat U(x)$，问题三使用等权准则得到综合方案。本题进一步考虑不同应用场景对三个指标的偏好差异，并寻找对权重变化不敏感的鲁棒设计。设计变量仍限定在附件2的离散水平：
\[
x_1\in\{0,0.10,0.15,0.20,0.30\},\quad
x_2\in\{3.0,3.5,4.0,4.5\},\quad
x_3\in\{2,4,6,8,10\}.
\]

\section{权重综合评价模型}
由于 $R,P,U$ 的数值范围不同，先采用附件2样本极值进行极差归一化：
\begin{equation}
z_j(x)=\frac{\hat y_j(x)-y_j^{\min}}{y_j^{\max}-y_j^{\min}},\qquad (y_1,y_2,y_3)=(R,P,U).
\end{equation}
给定权重 $w=(w_R,w_P,w_U)$，满足 $w_R+w_P+w_U=1$ 且 $w_j\ge0$，综合得分为
\begin{equation}
S_w(x)=w_Rz_R(x)+w_Pz_P(x)+w_Uz_U(x).
\end{equation}
每个应用场景下选择 $S_w$ 最小的离散组合。

设置七类代表性场景：均衡场景 $(1/3,1/3,1/3)$；低热阻芯片场景 $(0.6,0.2,0.2)$；低压降泵功耗场景 $(0.2,0.6,0.2)$；高均温可靠性场景 $(0.2,0.2,0.6)$；以及更极端的热阻、压降和均温优先场景 $(0.7,0.15,0.15)$、$(0.15,0.7,0.15)$、$(0.15,0.15,0.7)$。

\section{权重变化结果}
使用问题二的三个二次响应面，对100个离散组合逐一计算 $S_w$，结果如下。

\begin{table}[h]\centering\caption{不同应用场景下的最优方案}
\begin{tabular}{c|ccc|ccc|c}
\toprule
场景 & $w_R$ & $w_P$ & $w_U$ & $x_1$ & $x_2$ & $x_3$ & $(\hat R,\hat P,\hat U)$\\
\midrule
均衡 & .333 & .333 & .333 & .20 & 4.0 & 2 & (.755886,.083295,.784046)\\
低热阻芯片 & .60 & .20 & .20 & .20 & 3.0 & 8 & (.729009,.152900,.819339)\\
低压降泵功耗 & .20 & .60 & .20 & .15 & 4.5 & 2 & (.757983,.073021,.791587)\\
高均温可靠性 & .20 & .20 & .60 & .20 & 4.0 & 2 & (.755886,.083295,.784046)\\
热阻优先 & .70 & .15 & .15 & .20 & 3.0 & 8 & (.729009,.152900,.819339)\\
压降优先 & .15 & .70 & .15 & .15 & 4.5 & 2 & (.757983,.073021,.791587)\\
均温优先 & .15 & .15 & .70 & .20 & 4.0 & 2 & (.755886,.083295,.784046)\\
\bottomrule
\end{tabular}
\end{table}

结果表明，权重变化会引起方案切换，但切换集中在三类结构：热阻优先时采用较少歧管深度和较多针肋排数；压降优先时采用较大歧管深高比和较少排数；均衡或均温优先时回到 $(0.20,4.0,2)$。这说明三个设计变量分别承担换热强化、流阻降低和温度均匀化的作用。

\section{鲁棒设计准则}
仅选某一个权重下的最优方案可能在偏好改变后迅速失效。因此在权重集合 $\mathcal W$ 上采用 minimax 鲁棒准则：
\begin{equation}
x^{rob}=\arg\min_{x\in\Omega}\max_{w\in\mathcal W}S_w(x).
\end{equation}
为覆盖不同偏好，取 $w_R,w_P,w_U\ge0.1$、权重和为1的0.1间隔网格，共36组权重。对每个设计计算其36个得分，取最大值作为最坏情形得分，再选择最小者。

\section{鲁棒方案与优势}
计算得到鲁棒方案为
\begin{equation}
\boxed{x_1=0.20,\qquad x_2=3.5,\qquad x_3=6.}
\end{equation}
其代理模型预测值为
\[
(\hat R,\hat P,\hat U)=(0.740179,\ 0.123082,\ 0.797950).
\]
在36组权重下的最大综合得分为 $0.350780$，平均综合得分为 $0.319299$。该方案不是任何极端偏好下的单项冠军，却在热阻、压降和温度均匀性之间保持中等水平，因此权重改变时性能退化较慢。

鲁棒方案的物理解释为：针肋宽度比0.20处于强化换热和阻塞效应的折中点；歧管深高比3.5比4.5具有更强的换热能力，同时压降未达到3.0时的高水平；针肋排数6比2具有更低热阻，又避免8--10排导致的温度非均匀性明显上升。

\section{结论}
权重敏感性分析表明：热阻优先方案为 $(0.20,3.0,8)$，压降优先方案为 $(0.15,4.5,2)$，均衡和均温优先方案为 $(0.20,4.0,2)$。在允许三项指标权重均不低于0.1的偏好集合内，采用最坏情形综合得分最小的 minimax 准则，得到鲁棒方案 $(0.20,3.5,6)$。该方案对应用偏好变化不敏感，适合作为缺少明确权重或工况偏好可能变化时的推荐设计。

\end{document}
